% ==============================================================================
% Resumo em língua estrangeira (obrigatório, NBR 14724:2024)
% ==============================================================================
\begin{resumo}[Abstract]
\begin{instrucao}
\orientacao{Abstract}
Versão do resumo em inglês, com o mesmo conteúdo e as mesmas regras. Não use
tradução automática sem revisão cuidadosa. Termos técnicos devem seguir o uso
da área (por exemplo, \emph{multidimensional knapsack problem}, e não uma
tradução literal). Se usar ferramenta de IA para traduzir ou revisar, declare
o uso no apêndice sobre inteligência artificial.
\end{instrucao}

\begin{otherlanguage*}{english}
The multidimensional knapsack problem consists of selecting a subset of items
with maximum total value without exceeding the capacity of several resources at
the same time. It models decisions such as budget allocation among projects and
task selection on servers with processing and memory limits. Since the problem
is NP-hard, large instances are usually not solved to optimality in reasonable
time by integer linear programming solvers. This work proposes an iterated local
search heuristic that combines a greedy initial solution, based on an efficiency
ratio weighted by the constraints, with a swap-based local search and an
adaptive perturbation. The heuristic was implemented in C++ and compared with
the HiGHS solver and a genetic algorithm from the literature on 270 OR-Library
instances, with a time limit of 60 seconds per instance. The heuristic achieved
an average gap of 0.41\% from the best known solutions, against 1.87\% for the
solver in the same time, and outperformed the genetic algorithm in 82\% of the
instances. The results indicate that the adaptive perturbation is the component
that contributes most to performance, especially on instances with 500 items.

\keywords{multidimensional knapsack; iterated local search; heuristic;
combinatorial optimization}
\end{otherlanguage*}
\end{resumo}
